%!TEX TS-program = xelatex
%!TEX encoding = UTF-8 Unicode
% Awesome CV LaTeX Template for CV/Resume
%
% This template has been downloaded from:
% https://github.com/posquit0/Awesome-CV
%
% Customized for: Jaume Ivars Grimalt
% AI Engineer CV for data and AI consultancies. Same layout as resume/resume.tex.


%-------------------------------------------------------------------------------
% CONFIGURATIONS
%-------------------------------------------------------------------------------
\documentclass[11pt, a4paper]{../awesome-cv}
\usepackage{../cv-common}

%-------------------------------------------------------------------------------
%	PERSONAL INFORMATION
%-------------------------------------------------------------------------------
\name{Jaume}{Ivars Grimalt}
\position{AI Engineer}
% Location + relocation intent. No freelance / part-time availability.
\address{Currently in Valencia, Spain | EU citizen | Open to relocation to Zurich, Munich, London}

\mobile{+34 690 183 644}
\email{jaumeig95@gmail.com}
\homepage{jaumeivars.com}
\github{jaume2000}
\linkedin{jaume-ivars-grimalt}
\googlescholar{b5gMLtMAAAAJ}{Google Scholar}

\photo[circle,edge,left]{../PFP.jpeg}


%-------------------------------------------------------------------------------
\begin{document}

% Print the header with above personal information
\makecvheader[C]

% Print the footer with 3 arguments(<left>, <center>, <right>)
\makecvfooter
  {\today}
  {Jaume Ivars Grimalt~~~·~~~CV}
  {\thepage}


%-------------------------------------------------------------------------------
%	CONTENT
%-------------------------------------------------------------------------------
\cvsection{Profile}

\begin{cvparagraph}
I take models from the idea to production: data, training, deployment, backend and cloud. At Mycrospace I was the only engineer and built the whole stack, from a Master's thesis to a computer vision product used by 700 users in 35 countries.
\end{cvparagraph}

\cvsection{Technical Skills}
\begin{cvskills}
  \cvskill
    {ML/DL}
    {PyTorch (DDP, AMP, torch.compile), TensorFlow/TFLite, Hugging Face, timm, scikit-learn, OpenCV, pandas, ONNX}
  \cvskill
    {MLOps \& cloud}
    {GCP (Cloud Run, Cloud SQL, GCS), Hetzner/Dokploy, Docker, Redis/BullMQ, CI/CD, SLURM}
  \cvskill
    {Backend \& data}
    {Python, Flask, NestJS, PostgreSQL/Prisma, MongoDB}
  \cvskill
    {Generative AI}
    {LLMs via API, agents, diffusion (LDM, ControlNet, IP-Adapter)}
\end{cvskills}

\cvsection{Experience}
\begin{cventries}
\cvkeep
\cventry
    {Co-Founder \& CTO}
    {Mycrospace}
    {Valencia, Spain}
    {Jan 2024 – Oct 2026}
    {
\begin{cvitems}
\item {Computer vision product for microbiology labs. Master's thesis (2024) to company (Jan 2026). Sole engineer; built the whole stack. \textbf{700} users in 35 countries; tested in \textbf{5 labs}. 1st Prize IdeasUPV (2024).}
\item {\textbf{Production.} Event-based image pipeline: BullMQ/Redis queues, a Python Flask inference service, signed-URL uploads and SSE/WebSocket. Docker on GCP (Cloud Run, Cloud SQL, GCS), migrated to Hetzner/Dokploy to cut costs. NestJS/Prisma/PostgreSQL, MongoDB, Next.js, React Native/Expo and Stripe.}
\item {\textbf{MLOps.} \textbf{1,500}/\textbf{5,000} images annotated with microbiologists. Annotation versioning in MongoDB, k-fold CV, WBF ensembles, Cleanlab label audits and retraining on real client data. Vision models quantized and compiled with TFLite for edge.}
\item {\textbf{Models.} Own ConvNeXt detector (YOLO and DQ-DETR-inspired) for small objects and dense scenes; ~10–2000 colonies/image. \textbf{67 mAP50-95}; 92\% recall / 93\% precision. \textbf{12,500+} GPU-hours on Leonardo (two EuroHPC AI Factory allocations), \textbf{4x A100}, DDP/AMP/torch.compile on SLURM, checkpoint-resumed jobs.}
\end{cvitems}
    }
\cvkeep
\cventry
    {Research Engineer (full-time)}
    {MIALAB: Medical Imaging Analysis Laboratory, UPV}
    {Valencia, Spain}
    {May 2025 – Jan 2026}
    {
\begin{cvitems}
\item {Volumetric brain MRI segmentation with nnU-Net. Profiled volBrain (\textbf{700k+} volumes) from NIfTI ingest to voxel prediction for inference cost.}
\end{cvitems}
    }
\cvkeep
\cventry
    {Full-Stack Engineer (freelance)}
    {Tapstar}
    {Valencia, Spain}
    {Jan 2025 – May 2025}
    {
\begin{cvitems}
\item {Client delivery: QR generation, a back office and the frontend.}
\end{cvitems}
    }
\cvkeep
\cventry
    {Data Engineer}
    {Vesta-Z}
    {Valencia, Spain}
    {Jul 2023 – Feb 2025}
    {
\begin{cvitems}
\item {Cloud product features and a custom database schema so items could be cross-matched across providers. Delivered an LLM customer-support bot.}
\end{cvitems}
    }
\end{cventries}
\vspace{0.4mm}
\noindent{\fontsize{8pt}{1em}\bodyfont\bfseries\color{darktext} Earlier experience:}~~{\fontsize{8pt}{1em}\bodyfontlight\color{text} Neurocatching, Data Scientist Intern, 2023; Sciling, Backend Developer Intern, 2022.}\par

\cvsection{Generative AI}
\par\nopagebreak
\vspace{\acvSectionContentTopSkip}
\vspace{-2.5mm}
{\bodyfont\fontsize{7.6pt}{0.95em}\selectfont\raggedright
\noindent\textbullet~~Python pipeline of LLM-based agents that automates profile research and drafts outreach emails and messages.\par
\vspace{0.15mm}
\noindent\textbullet~~Prototype, not deployed: synthetic data augmentation with diffusion (LDM, ControlNet and IP-Adapter).\par
}

%-------------------------------------------------------------------------------
% SECTION TITLE
%-------------------------------------------------------------------------------
\cvsection{Education}
\nopagebreak
\begin{cventries}
\cvkeep
\cvdegree
    {UPV · GPA 8.7/10 · Thesis Honours (became Mycrospace)}
    {M.Sc. Artificial Intelligence, Pattern Recognition \& Digital Imaging}
    {Valencia, Spain}
    {2023 – 2024}
\cvkeep
\cvdegree
    {CDTM (TUM/LMU) · 16 selected}
    {Technology Management}
    {Valencia / Munich}
    {2024}
\cvkeep
\cvdegree
    {UPV · GPA 8.3/10 · Honours in ML and Statistics}
    {B.Sc. Computer Engineering}
    {Valencia, Spain}
    {2019 – 2023}
\end{cventries}

\vspace{0.4mm}
\noindent{\fontsize{8pt}{1em}\bodyfont\bfseries\color{darktext} Teaching:}~~{\fontsize{8pt}{1em}\bodyfontlight\color{text} Big Data module, CEI, Valencia, Oct 2026 – Dec 2026.}\par

\cvsection{Publications}
\par\nopagebreak
\vspace{\acvSectionContentTopSkip}
\vspace{-2.5mm}
{\bodyfont\fontsize{7.6pt}{0.95em}\selectfont\raggedright
\noindent\textbullet~~\textit{Auto-N2N: Self-supervised MRI Denoising using Synthetic Pair Generation and Mixture of Experts.}
S.\ Morell-Ortega et al., incl.\ \textbf{J.\ Ivars Grimalt}. Preprint, 2026. \href{https://openreview.net/forum?noteId=1f93N1n8RP}{OpenReview}\par
\vspace{0.15mm}
\noindent\textbullet~~\textit{Detecting Cytopathic Effect in Cell Culture with Computer Vision.}
W.\ Ragazzo, \textbf{J.\ Ivars Grimalt}. Under review.\par
\vspace{0.15mm}
\noindent\textbullet~~\textit{What Extra Euler Steps Reveal About Residual Networks: Feature Dynamics in Weight-Tied and Untied Stages.} Manuscript in preparation.\par
\vspace{0.6mm}
\noindent{\fontsize{8pt}{1em}\bodyfont\bfseries\color{darktext} Languages:}~~Spanish (native), Catalan (native), English (C1)\par
}

%-------------------------------------------------------------------------------
\end{document}
